% SPDX-License-Identifier: Apache-2.0
\section{\code{txhdl::foreign}}
\label{sec:r-foreign}

A module written in Verilog as a unit. Verilator compiles the module
into a C++ model, and \code{//tools/vshim} writes two things from the
module's port list: a C shim over the model, and the Rust that makes
a \code{Unit} of it. The shim is four calls and a handle: settle the
combinational logic on the inputs as they stand, pulse the clock,
set an input port by index, read a port by index, values as four
words of thirty-two bits, so a port up to 128 bits wide crosses.
\code{Model} holds the handle and the shim's functions and frees the
model when dropped. The generated unit has an \code{In} or
\code{Out} per port and an \code{Rx} or \code{Tx} per \code{x\_data},
\code{x\_valid}, \code{x\_ready} trio, and its step, once per rising
edge, is: the inputs in as the edge left them, a channel's head and
its offer, and the room the runtime has for the module's sends; a
settle; the transfers the module agrees to, a take where it says
ready and an offer stands, a send where it says valid and there is
room; the edge; and the outputs out. So the module is a registered
unit, its outputs seen a step later, and goes first in the join
order as the timer of the core does. What it cannot do is lower: the
lowering carries it as nothing, and a netlist that holds one is the
board top's business, by hand, as the core's is.
\code{verilog\_unit()} in \code{lib/foreign.bzl} builds the model,
the shim and the crate from the one Verilog file.

A VHDL entity goes the same way with nvc as the engine, run as a
child process, since nvc is a program and not a library. The tool
writes a testbench around the entity that reads a line per step
from its standard input, a command and the inputs as hex, applies
them, settles or pulses the clock, and writes the outputs back as
one line; the build analyses and elaborates the two and makes the
script that runs nvc on them, an executable of the build's own,
since a test target cannot be depended on; and \code{Cosim} starts
the script in the runfiles, where its paths hold, and speaks the
lines, a settle and an edge each one exchange. The same generated
unit drives either engine, since both answer to set, get, settle
and edge; a step costs a pipe's round trip under nvc, some tens of
microseconds, and nothing under Verilator. \code{vhdl\_unit()}
builds the testbench, the script and the crate. The examples
document~\cite{examples} co-runs a unit with its own Verilog and
its own VHDL at once.

\lstinputlisting[caption={\code{lib/src/foreign.rs}.},label={lst:r-foreign}]{lib/src/foreign.rs}
